\documentclass[11pt]{amsart}
\usepackage[margin=1.1in]{geometry}
\usepackage{amsmath,amssymb,amsthm}
\usepackage[colorlinks=true,linkcolor=blue,citecolor=blue,urlcolor=blue]{hyperref}

\newtheorem{theorem}{Theorem}
\newtheorem{lemma}[theorem]{Lemma}
\newtheorem{proposition}[theorem]{Proposition}
\newtheorem{corollary}[theorem]{Corollary}
\theoremstyle{remark}
\newtheorem{remark}[theorem]{Remark}

\newcommand{\lean}[1]{\texttt{#1}}
\setlength{\emergencystretch}{3em}

\begin{document}

\title[Divisibility of a factorial ratio]{Divisibility of the factorial ratio $(8n)!\,n!/\bigl((4n)!\,(3n)!\,(2n)!\bigr)$: proofs of conjectures of Bala}

\author{Ansh Mishra}
\author{Aryan Senthilkumar}

\date{October 2026}

\begin{abstract}
Let $a(n)=(8n)!\,n!/((4n)!\,(3n)!\,(2n)!)$ (OEIS A211420).
Peter Bala conjectured that $(n+1)\mid35\,a(n)$, $(2n+1)\mid3\,a(n)$, $(3n+1)\mid5\,a(n)$ and $(8n-1)\mid a(n)$ for all $n$.
He also conjectured that the products $(kn+1)\cdots(kn+r)$ for $k\in\{1,2,3\}$ and $(8n-1)(8n-3)\cdots(8n-2r-1)$ divide fixed multiples of $a(n)$.
We prove all of these.
The divisibilities $(n+1)\mid35a(n)$, $(2n+1)\mid3a(n)$ and $(8n-1)\mid a(n)$ appear to be new, and the constants $35$, $3$ and $5$ are optimal.
The proofs evaluate Legendre's formula through a single step function of period $24$.
At the primes $2$ and $3$ we use a short ``escape'' lemma.
All four statements, as formalized in Google DeepMind's Formal Conjectures repository, are proved in Lean~4 with Mathlib.
\end{abstract}

\maketitle

\section{Introduction}

The numbers $a(n)=\frac{(8n)!\,n!}{(4n)!\,(3n)!\,(2n)!}=1,140,60060,29745716,\dots$ form OEIS A211420.
They are integers by Landau's criterion; they also belong to a family studied by Bober \cite{Bober}.
In 2025 Bala recorded the following conjectures in the OEIS entry \cite{OEIS}:
\begin{enumerate}
\item[(C1)] $(n+1)\mid35\,a(n)$, $(2n+1)\mid3\,a(n)$ and $(3n+1)\mid5\,a(n)$ for all $n\ge0$;
\item[(C2)] for $k\in\{1,2,3\}$ and $r\ge1$ there is $C(k,r)>0$ with $(kn+1)(kn+2)\cdots(kn+r)\mid C(k,r)\,a(n)$ for all $n$;
\item[(C3)] $(8n-1)\mid a(n)$ for all $n$;
\item[(C4)] for $r\ge0$ there is $K(r)>0$ with $\prod_{i=0}^{r}(8n-2i-1)\mid K(r)\,a(n)$ for all $n$.
\end{enumerate}
They are stated in the Formal Conjectures repository \cite{FC} as \lean{OeisA211420.conjecture1}, \dots, \lean{conjecture4}.

\begin{theorem}\label{thm:main}
(C1)--(C4) hold.
In (C1) the constants $35,3,5$ cannot be decreased.
One can take $C(k,r)=((24r)!)^r$ and $K(r)=((6r+3)!)^{r+1}$.
\end{theorem}

\emph{Prior work.}
The divisibility $(3n+1)\mid5a(n)$ is the case $k=4$ of Sun's theorem $\binom{2n}{n}\mid(k+1)'\,C^{(k-1)}_n\binom{2kn}{kn}$ \cite{Sun}, because $a(n)=\binom{8n}{4n}\binom{4n}{n}/\binom{2n}{n}$.
Yang \cite{Yang} proved $(2n+1)(2n+3)\mid99\,a(n)$.
Statement (C2), with the explicit constant $\mathrm{lcm}(1,\dots,8r-1)^r$, was proved and formalized in Lean in \cite{Danny}.
Statement (C4), with $K(r)=(8r)!$, was proved by Adamczewski \cite[Appendix~A]{Adam}.
For $r=0$ this gives only $(8n-1)\mid8!\,a(n)$.
We found no earlier proof of $(n+1)\mid35a(n)$, $(2n+1)\mid3a(n)$ or $(8n-1)\mid a(n)$.

\section{Legendre's formula as a residue function}

For a prime power $q=p^j$ put $T=\lfloor24n/q\rfloor$.
For $c\in\{1,2,3,4,8\}$ we have $c\mid24$, and therefore $\lfloor cn/q\rfloor=\lfloor T/(24/c)\rfloor$.
So Legendre's formula gives
\begin{equation}\label{eq:vp}
v_p(a(n))=\sum_{j\ge1}\varphi\bigl(\lfloor24n/p^j\rfloor\bigr),\qquad \varphi(T)=\Bigl\lfloor\frac T3\Bigr\rfloor+\Bigl\lfloor\frac T{24}\Bigr\rfloor-\Bigl\lfloor\frac T6\Bigr\rfloor-\Bigl\lfloor\frac T8\Bigr\rfloor-\Bigl\lfloor\frac T{12}\Bigr\rfloor .
\end{equation}
Since $8+1-4-3-2=0$, $\varphi$ has period $24$.
By inspection of $T=0,\dots,23$,
\begin{equation}\label{eq:G}
\varphi(T)=1\iff T\bmod24\in G:=\{3,4,5,6,7,9,10,11,15,21,22,23\},
\end{equation}
and $\varphi(T)=0$ otherwise.
In particular $\varphi\ge0$, which is the integrality of $a(n)$.
To prove $t\mid c\,a(n)$ it suffices, for every prime $p$ and every $j\ge1$ with $p^j\mid t$, that $p^j\mid c$ or $\varphi(\lfloor24n/p^j\rfloor)=1$.
We call this the \emph{termwise} criterion.
Write $T_j=\lfloor24n/p^j\rfloor$, so that $T_{j+1}=\lfloor T_j/p\rfloor$.

\begin{lemma}[location]\label{lem:loc}
Let $q\ge2$ and let $m$ denote the cofactor in each case.
\begin{enumerate}
\item[(L1)] If $q\mid n+1$ and $q\ge8$, then $T\bmod24\in\{21,22,23\}$.
\item[(L2)] If $q\mid2n+1$ and $q\ge4$, then $T\bmod24\in\{9,10,11\}$.
\item[(L3)] If $q\mid3n+1$, $q\ge4$ and $q\ne5$, then $T\bmod 24\in\{6,7,15,23\}$.
\item[(L4)] If $q\mid8n-1$ and $q\ge4$, then $T=3m$ with $m$ odd, so $T\bmod24\in\{3,9,15,21\}$.
\end{enumerate}
In all four cases $\varphi(T)=1$.
\end{lemma}

\begin{proof}
(L1): $n+1=qm$ gives $24n=24mq-24$ with $24\le3q$, so $24m-3\le T\le24m-1$.
(L2): $m$ is odd and $24n=12mq-12$ with $12\le3q$, so $12m-3\le T\le12m-1$.
(L4): $8n-1=qm$ gives $24n=3mq+3$ with $3<q$, so $T=3m$.
(L3): for $q\ge8$ we get $T=8m-1$, whose residue is $7$, $15$ or $23$. For $q=4$, $n\equiv1\pmod4$ gives $T=6n\equiv6$. For $q=7$, $n\equiv2\pmod7$ gives $T\equiv6$. The case $q=6$ cannot occur.
\end{proof}

\begin{lemma}[escape]\label{lem:escape}
Let $p=3$ and suppose $T_i\bmod24\ge9$; or let $p=2$ and suppose $T_i\bmod24\ge6$.
Then $\varphi(T_j)=1$ for some $j>i$.
\end{lemma}

\begin{proof}
Let $p=3$ and suppose $\varphi(T_j)=0$ for all $j>i$.
We show by downward induction that $T_j\bmod24\le8$ for all $j\ge i$; the start is $T_j=0$ for large $j$.
If $S=T_{j+1}$ has $\varphi(S)=0$ and $S\bmod24\le8$, then $S\bmod24\in\{0,1,2,8\}$ by \eqref{eq:G}, so $S\bmod8\le2$.
Hence $T_j\bmod 24=3(S\bmod8)+(T_j\bmod3)\le8$.
This contradicts $T_i\bmod24\ge9$.
For $p=2$ the same argument works with $S\bmod24\in\{0,1,2\}$ and $T_j\bmod24=2(S\bmod12)+(T_j\bmod2)\le5$.
\end{proof}

\section{Proofs}

\begin{proof}[Proof of (C3)]
At $n=0$ the divisor is $-1$.
Let $n\ge1$; then $t=8n-1$ is odd.
For $p\ge5$ every $p^j$ is at least $5$, so the termwise criterion holds by (L4).
For $p=3$ let $e=v_3(t)\ge1$ and $t=3^em$.
The indices $j=2,\dots,e$ contribute $e-1$ by (L4).
Moreover $24n=3^{e+1}m+3$ gives $T_{e+1}=m$.
From $3^em\equiv7\pmod8$ we get $m\equiv5$ or $7\pmod8$, and $3\nmid m$.
So $m\bmod24\in\{5,7,13,23\}$.
Either $\varphi(m)=1$, or $m\equiv13$, in which case Lemma~\ref{lem:escape} gives an index $j>e+1$ with $\varphi(T_j)=1$.
Hence $v_3(a(n))\ge e$.
\end{proof}

\begin{proof}[Proof of (C1)]
\emph{$(2n+1)\mid3a(n)$.}
Here 2 does not divide $2n+1$, so let $q=p^j\mid2n+1$ with $p$ odd.
If $q\ge4$, (L2) applies; if $q=3$, then $q\mid3$.

\emph{$(3n+1)\mid5a(n)$.}
For odd $p$ the criterion holds termwise by (L3), since $3\nmid3n+1$ and $5\mid5$.
For $p=2$ let $e=v_2(3n+1)$ and $3n+1=2^eu$.
The indices $j=2,\dots,e$ contribute $e-1$ by (L3).
If $e=1$, then $n\equiv3\pmod4$ and $T_2=6n\equiv18\pmod{24}$.
If $e\ge2$, then $T_{e+1}=4u-1$, whose residue is $3$, $11$ or $19$.
In each case $\varphi(T_{e+1})=1$, or $T_{e+1}\bmod24\ge6$ and Lemma~\ref{lem:escape} supplies the missing index.

\emph{$(n+1)\mid35a(n)$.}
For $p\notin\{2,3\}$ the criterion holds termwise: by (L1) if $q\ge8$, while $q<8$ forces $q=p\in\{5,7\}$, and $p\mid35$.
For $p=3$ let $e=v_3(n+1)$ and $n+1=3^em$.
The indices $j=2,\dots,e$ contribute $e-1$ by (L1).
Also $T_{e+1}\in[8m-3,8m-1]$, with residue in $\{5,6,7,13,14,15,21,22,23\}$.
So $\varphi(T_{e+1})=1$, or the residue is $13$ or $14$ and Lemma~\ref{lem:escape} applies.
For $p=2$ let $e=v_2(n+1)$ and $n+1=2^eu$.
The indices $j=3,\dots,e$ contribute $e-2$ by (L1).
The value $T_{e+1}$ equals $12u-6$, $12u-3$, or lies in $[12u-2,12u-1]$, according as $e=1$, $e=2$ or $e\ge3$.
Its residue is $6$, $9$, $10$ or $11$, so $\varphi(T_{e+1})=1$ and $T_{e+1}\bmod24\ge6$.
Lemma~\ref{lem:escape} then gives a further index, for a total of at least $e$.

\emph{Optimality.}
$7\mid n+1$ but $a(6)\equiv1\pmod7$; $5\mid n+1$ but $a(9)\equiv1\pmod5$; $a(1)=140\not\equiv0\pmod3$; and $3n+1=10$ while $a(3)=29745716\not\equiv0\pmod5$.
\end{proof}

\begin{proof}[Proof of (C2) and (C4)]
For (C2), let $P=(kn+1)\cdots(kn+r)$.
Then $v_p(P)=\sum_jc_j$, where $c_j$ is the number of multiples of $q=p^j$ in $(kn,kn+r]$.
If $q\le24r$, then $c_j\le r\le r\lfloor24r/q\rfloor$, which is covered by $v_p(((24r)!)^r)$.
If $q>24r$, then $c_j\le1$.
When $c_j=1$ the multiple is $kn+i$ with $24i<q$; writing $D=24/k$, we get $T_j=Dm-1$ and hence $\varphi(T_j)=1$.

For (C4), let $\ell_i=|8n-2i-1|$, which is odd and nonzero.
Prime powers $q\le6r+3$ are covered by $v_p(((6r+3)!)^{r+1})$.
If $q>6r+3$ and $q\mid\ell_i$, then $8n>2i+1$ and $8n=qm+2i+1$ with $3(2i+1)<q$.
So $T_j=3m$ with $m$ odd and $\varphi(T_j)=1$.
Since $m$ determines $i$, at most one $\ell_i$ is divisible by $q$.
\end{proof}

\section{Formalization and checks}

All four statements of \cite{FC} are proved in Lean~4 with Mathlib, against the unmodified statement file.
The theorems \lean{OeisA211420.conjecture1\_solved}, \dots, \lean{conjecture4\_solved} have the types of the repository statements; this is checked by \lean{rfl}.
They depend only on the axioms \lean{propext}, \lean{Classical.choice} and \lean{Quot.sound}.
The formal proofs follow Sections~2--3, with Legendre's formula taken from Mathlib.
We also checked (C1) and (C3) numerically for all $n\le200\,000$ via \eqref{eq:vp}, and with exact integers for $n<400$.
(C2) and (C4) were checked for $r\le6$ and $n\le3000$.

\begin{thebibliography}{9}

\bibitem{Adam}
T.~Adamczewski, \emph{OEIS Open: How many conjectures can language models turn into theorems?}, arXiv:2608.11941, 2026.

\bibitem{Bober}
J.~W.~Bober, \emph{Factorial ratios, hypergeometric series, and a family of step functions}, J.\ London Math.\ Soc.\ \textbf{79} (2009), 422--444.

\bibitem{Danny}
DannyExperiments, \emph{A211420 fixed-divisor theorem}, \url{https://github.com/DannyExperiments/a211420-fixed-divisor}, Zenodo, doi:10.5281/zenodo.21762443, 2026.

\bibitem{FC}
Google DeepMind, \emph{Formal Conjectures}, \url{https://github.com/google-deepmind/formal-conjectures}; the statements are in \texttt{FormalConjectures/OEIS/211420.lean}, pull request 6426 (2026).

\bibitem{OEIS}
OEIS Foundation Inc., \emph{The On-Line Encyclopedia of Integer Sequences}, \url{https://oeis.org/A211420}.

\bibitem{Sun}
Z.-W.~Sun, \emph{On divisibility of binomial coefficients}, J.\ Aust.\ Math.\ Soc.\ \textbf{93} (2012), 189--201.

\bibitem{Yang}
Q.~Yang (as cited in \cite{Danny}), Electron.\ J.\ Combin.\ \textbf{22} (2015), \#P1.9.

\end{thebibliography}

\end{document}
